\begin{table*}[!t]
\centering
\small
\renewcommand\tabcolsep{1.0pt} % column space
\begin{tabular}{cp{11.5cm}}
    \toprule
    \textbf{Evaluator} & \textbf{Prompt} \\
    \midrule
    \multirow{8}{*}{\parbox{4cm}{\centering LLM-based Evaluator \\(0-0.5-1 setting)}}
    &  I am testing a model's performance on open-ended questions. I want you to help me in checking to see if the candidate answer has the same meaning as the reference answer for a given question. If you think the reference answer and the candidate answer have the same meaning, respond \{``selection'': ``1''\}; otherwise, respond by \{``selection'': ``0''\}. If you think the candidate is partially correct, respond by \{``selection'': ``0.5''\}. If the answer is ``I don't know,'' rate it to 0.

    <QUESTION>
    \{question\}
    </QUESTION>

    <REFERENCE>
    \{reference\}
    </REFERENCE>

    <CANDIDATE>
    \{candidate\}
    </CANDIDATE>
\\
\midrule

\multirow{28}{*}{\parbox{4cm}{\centering LLM-based Evaluator \\(1-2-3-4-5 setting)}}
    &  I am testing a model's performance on open-ended questions. I want you to help me in checking to see if the candidate answer has the same meaning as the reference answer for a given question. 
    
    Rate the candidate answer from 1, 2, 3, 4, and 5, where 1 means the candidate is the least similar to the reference answer and 5 means the candidate matches to the reference answer perfectly. 
    Respond by \{``selection'': ``''\}. If the candidate answer is ``I don't know,'' rate it to 1.

    Here's some examples you can consider:
    
    Question: Why transformer is better than RNN?
    
    Reference: Parallel computation
    
    Candidate: Computation
    
    Rating: 3

    Question: What's the major advantage of using ALiBi positional embedding?
    
    Reference: Effectively handle sequences of varying lengths, particularly beneficial for very long sequences
    
    Candidate: It has more freedom to handle input
    
    Rating: 2

    Question: What's the model's performance on GSK8K dataset? 
    
    Refernece: 65.65\%
    
    Candidate: 44.56\%
    
    Rating: 1

    Question: What specific method does this paper propose to solve LLM searching problem?
    
    Reference: MCTS
    
    Candiate: Monte Carlo Tree Search is proposed in this paper to solve searching when using decomposed prompting method.
    
    Rating: 5

    Question: How does the performance change when we switch from CoT to ToT in prompting? 
    
    Reference: Accuracy from 23.50\% to 32.87\%
    
    Candidate: slightly increase
    
    Rating: 4

    <QUESTION>
    \{question\}
    </QUESTION>

    <REFERENCE>
    \{reference\}
    </REFERENCE>

    <CANDIDATE>
    \{candidate\}
    </CANDIDATE>
\\
    \bottomrule
\end{tabular}
\caption{Prompts used to evaluate answers generated by LLMs.}
\label{tab:gpt_based_evaluator_app}
\end{table*}